\chapter{Конструкторский раздел}

\section{Основной сценарий протокола}

Диаграмма последовательности отражает порядок взаимодействия между пользователем (клиентским приложением) и сервером в рамках разработанного протокола совместного редактирования текстовых заметок. Взаимодействие реализуется по модели клиент–сервер и включает этапы аутентификации, управления заметками, синхронизации данных и разрешения конфликтов при параллельном редактировании.

На рисунке 2.2 представлена диаграмма последовательности работы протокола совместного редактирования заметок.

\begin{figure}[H]
    \centering
    \includegraphics[width=0.94\textwidth]{../img/CW_network.drawio.pdf}
    \caption{Диаграмма последовательности работы протокола}
    \label{fig:interaction}
\end{figure}

Процесс взаимодействия начинается с регистрации пользователя в системе. Клиент отправляет на сервер запрос регистрации, содержащий учётные данные пользователя. Сервер выполняет проверку корректности введённых данных и возвращает клиенту результат регистрации. В случае успешной регистрации пользователь получает возможность пройти процедуру авторизации.

После авторизации пользователь может создавать новые заметки. Для этого клиент формирует запрос на создание заметки, который передаётся на сервер. Сервер проверяет корректность запроса, создаёт новую заметку и возвращает клиенту результат операции, включая идентификатор заметки и её начальную версию.

Далее пользователь может открыть существующую заметку. Клиент отправляет запрос на открытие заметки, после чего сервер возвращает актуальный текст заметки и номер её версии. Полученные данные используются клиентом для отображения содержимого заметки в пользовательском интерфейсе.

В процессе редактирования пользователь вносит изменения в текст заметки и отправляет их на сервер. Сервер сохраняет изменения и увеличивает номер версии заметки. При отсутствии конфликтов сервер подтверждает успешное сохранение изменений.

В случае совместного доступа к заметке сервер рассылает уведомления другим пользователям о наличии общей заметки и изменениях в её содержимом. Таким образом обеспечивается согласованность данных между всеми участниками редактирования.

Если во время сохранения изменений сервер обнаруживает конфликт версий (например, если другой пользователь успел изменить заметку ранее), запускается механизм разрешения конфликтов. На первом этапе сервер выполняет попытку автоматического разрешения конфликта, анализируя внесённые изменения. Если изменения пользователей не пересекаются или могут быть объединены без потери данных (например, изменения затрагивают разные фрагменты текста), сервер формирует согласованную версию заметки автоматически.

В случае успешного автоматического разрешения конфликта сервер сохраняет новую версию заметки и рассылает обновлённый текст всем пользователям, имеющим доступ к данной заметке. Клиенты получают актуальную версию заметки без необходимости дополнительного вмешательства пользователя.

Если автоматическое разрешение конфликта невозможно, сервер уведомляет пользователя о наличии конфликта и предлагает варианты дальнейших действий. Пользователю предоставляется возможность выбрать способ разрешения конфликта.

При выборе автоматической синхронизации клиент запрашивает у сервера актуальную серверную версию заметки. Сервер возвращает текущий текст и номер версии заметки, после чего пользователь может принять серверную версию и продолжить работу.

Если пользователь выбирает ручное разрешение конфликта, клиент формирует новую согласованную версию заметки и отправляет её на сервер на согласование. Сервер, в свою очередь, пересылает запрос владельцу заметки для подтверждения изменений.

Владелец заметки может либо принять изменения, либо отклонить их. В случае принятия сервер сохраняет новую версию заметки и рассылает обновлённый текст всем пользователям, имеющим доступ к данной заметке. При отказе сервер уведомляет пользователя, инициировавшего согласование, об отклонении изменений.

Таким образом, разработанный протокол обеспечивает корректную синхронизацию данных, поддержку совместного редактирования, автоматическое и ручное разрешение конфликтов, а также детерминированное согласование изменений при параллельной работе нескольких пользователей с одной заметкой.

\section{Участники совместного редактирования заметок и их функции}

В системе совместного редактирования текстовых заметок задействованы несколько ключевых участников, выполняющих определённые функции для обеспечения корректного обмена данными, синхронизации версий документов и разрешения конфликтов при параллельном редактировании. На рисунке 2.1 представлена схема взаимодействия основных компонентов системы.

Система построена по клиент–серверной архитектуре, в рамках которой клиенты инициируют операции редактирования и управления заметками, а сервер отвечает за централизованную обработку запросов, хранение данных и согласование изменений между пользователями.

\begin{figure}[H]
	\centering
	{\includegraphics[width = \textwidth]{../img/client_server_interaction.jpg}}
	\caption{Участники системы совместного редактирования заметок}
	\label{fig:page-03}
\end{figure}

Клиент представляет собой пользовательское приложение, предназначенное для работы с текстовыми заметками. Клиент обеспечивает взаимодействие пользователя с системой, формирует запросы протокола и обрабатывает ответы сервера.

Основные функции клиента:

\begin{itemize}
    \item аутентификация и регистрация пользователя — формирование и отправка запросов AUTH и REGISTRATION для получения доступа к системе;
    \item получение списка доступных заметок — запрос GET\_NOTES, позволяющий получить перечень личных и совместных заметок;
    \item создание и открытие заметок — операции CREATE\_NOTE и OPEN\_NOTE;
    \item редактирование текста заметки — отправка операций изменения содержимого с использованием UPDATE\_TEXT;
    \item синхронизация состояния документа — запрос SYNC при рассинхронизации версий;
    \item совместный доступ к заметкам — инициирование операции SHARE\_NOTE и получение уведомлений SHARE\_NOTE\_NOTIFY;
    \item участие в разрешении конфликтов — отправка запросов APPROVE\_MERGE и обработка ответов OWNER\_APPROVE\_MERGE, SERVER\_APPROVE\_MERGE;
    \item обработка обновлённого текста — применение изменений, полученных в сообщении UPDATE\_TEXT\_MERGED.
\end{itemize}

Клиент также выполняет функции шифрования исходящих сообщений и дешифрования данных, получаемых от сервера.

Сервер является центральным компонентом системы и отвечает за приём, обработку и хранение данных, а также за координацию совместного редактирования между несколькими пользователями.

Основные функции сервера:

\begin{itemize}
    \item приём и обработка клиентских запросов — обработка сообщений всех типов, определённых в протоколе;
    \item аутентификация и регистрация пользователей — проверка учётных данных и управление идентификаторами пользователей (AUTH, REGISTRATION);
    \item управление заметками — создание, открытие и предоставление списка заметок (CREATE\_NOTE, OPEN\_NOTE, GET\_NOTES);
    \item контроль версий документов — отслеживание актуальных версий заметок и обработка запросов синхронизации (SYNC);
    \item обработка операций редактирования — применение изменений, полученных от клиентов (UPDATE\_TEXT);
    \item распространение изменений — уведомление всех участников совместного редактирования о новых версиях текста (UPDATE\_TEXT\_MERGED);
    \item разрешение конфликтов — координация процессов согласования изменений между пользователями (APPROVE\_MERGE, OWNER\_APPROVE\_MERGE, SERVER\_APPROVE\_MERGE);
    \item обеспечение целостности и безопасности данных — проверка корректности сообщений, дешифрование и валидация данных;
    \item уведомление клиентов — отправка ответов и событийных сообщений клиентским приложениям.
\end{itemize}

Таким образом, сервер обеспечивает согласованное состояние текстовых заметок у всех участников системы и выступает в роли центра принятия решений при возникновении конфликтных ситуаций.


\section{Спецификации для структуры сообщений протокола совместного редактирования заметок}

Проектирование структуры сообщений протокола совместного редактирования текстовых заметок требует соблюдения ряда критериев, обеспечивающих корректную интерпретацию, обработку и синхронизацию данных между участниками системы. Ниже представлены ключевые аспекты, которые необходимо учитывать при разработке архитектуры протокола взаимодействия клиентов и сервера.

Структура сообщений протокола должна обладать достаточной степенью гибкости, позволяющей расширять набор поддерживаемых операций и типов данных без существенных изменений базовой логики системы. Это особенно важно в условиях совместного редактирования, где в дальнейшем могут быть добавлены новые механизмы синхронизации, разрешения конфликтов или управления доступом.

Передаваемые сообщения должны содержать исключительно релевантную информацию, необходимую для выполнения конкретной операции. Исключение избыточных или дублирующих данных позволяет снизить сетевой трафик, уменьшить задержки передачи и повысить общую производительность системы, что критично при одновременной работе нескольких пользователей с одним документом.

Каждое сообщение протокола должно включать как полезную нагрузку, так и служебные метаданные. К таким метаданным относятся тип операции, идентификатор пользователя, идентификатор заметки, версия документа, а также длина сообщения. Наличие этих атрибутов обеспечивает корректную маршрутизацию сообщений, контроль целостности данных и возможность проверки актуальности изменений при синхронизации.

Особое внимание при проектировании структуры сообщений уделяется поддержке механизмов версионности. Информация о версии заметки позволяет серверу выявлять конфликты при одновременном редактировании текста несколькими пользователями и инициировать процедуры их разрешения, такие как откат изменений, объединение версий или запрос подтверждения у владельца заметки.

Элементы данных в сообщениях протокола должны быть взаимосвязаны таким образом, чтобы сохранять логическую целостность состояния заметки на всех клиентах системы. Это обеспечивает непротиворечивость представления данных, упрощает последующую обработку сообщений на стороне сервера и гарантирует согласованность содержимого документа у всех участников совместного редактирования.

\section{Описание структур сообщений}

Разработанный протокол совместного редактирования текстовых заметок использует бинарный формат передачи данных с байтовой адресацией, обеспечивающий однозначность интерпретации сообщений и эффективность обработки на стороне клиента и сервера.

В рамках протокола применяется неоптимизированное по битам бинарное представление, при котором элементы сообщения кодируются с использованием целого числа байтов. Такой подход выбран в целях упрощения реализации, повышения читаемости структуры сообщений и снижения вероятности ошибок при сериализации и десериализации данных.

Использование бинарного формата позволяет минимизировать накладные расходы по сравнению с текстовыми протоколами и упростить реализацию механизмов синхронизации и контроля версий при совместном редактировании.

Каждый передаваемый пакет имеет фиксированную логическую структуру и состоит из двух основных частей:

\begin{itemize}
    \item заголовок — 1 байт, содержащий код операции протокола, однозначно определяющий тип передаваемого сообщения (запрос или ответ) и формат его последующей обработки;
    \item полезная нагрузка — переменной длины, формируемая в соответствии с типом операции и содержащая данные запроса или ответа, такие как идентификаторы пользователей и заметок, версии документа, текстовые изменения или служебную информацию, необходимую для синхронизации и разрешения конфликтов.
\end{itemize}

Формат полезной нагрузки зависит от значения кода операции и определяется спецификацией протокола для каждого типа сообщения. Такой подход обеспечивает гибкость протокола, позволяя расширять набор поддерживаемых операций без изменения базовой структуры пакета.

\begin{table}[H]
    \centering
    \caption{Атрибутивный состав запроса авторизации пользователя}
    \begin{tabular}{|p{4cm}|p{2cm}|p{3cm}|p{6cm}|}
    \hline
    \textbf{Параметр} & \textbf{Размер} & \textbf{Тип данных} & \textbf{Описание} \\ \hline
    Код операции & 1 байт & uint8\_t & Уникальный код операции (0x01 (AUTH) — запрос на авторизацию) \\ \hline
    Длина логина & 1 байт & uint8\_t & Длина логина пользователя \\ \hline
    Длина пароля & 1 байт & uint8\_t & Длина пароля пользователя \\ \hline
    Логин & L байт (\(L \le 255\)) & byte[L] (UTF-8) & Логин пользователя, длина определяется значением поля "Длина логина" \\ \hline
    Пароль & L байт (\(L \le 255\)) & byte[L] (UTF-8) & Пароль пользователя, длина определяется значением поля "Длина пароля" \\ \hline
    \end{tabular}
    \label{tab:auth_request}
\end{table}

\begin{table}[H]
    \centering
    \caption{Атрибутивный состав ответа сервера на запрос авторизации пользователя}
    \renewcommand{\arraystretch}{1.2}
    \begin{tabular}{|p{4cm}|p{2.2cm}|p{3cm}|p{6cm}|}
    \hline
    \textbf{Параметр} & \textbf{Размер} & \textbf{Тип данных} & \textbf{Описание} \\ \hline
    Код операции & 1 байт & uint8\_t & Уникальный код операции (0x01 (AUTH) — запрос на авторизацию) \\ \hline
    Статус выполнения & 1 байт & uint8\_t & 0 — успех, 1 — ошибка \\ \hline
    ID пользователя & 1 байт & uint8\_t & в случае status == 1, передается 0, в остальных случаях нумерация пользователей начинается 0 \\ \hline
    \end{tabular}
    \label{tab:auth_response}
\end{table}


\begin{table}[H]
    \centering
    \caption{Атрибутивный состав запроса регистрации нового пользователя}
    \begin{tabular}{|p{4cm}|p{2cm}|p{3cm}|p{6cm}|}
    \hline
    \textbf{Параметр} & \textbf{Размер} & \textbf{Тип данных} & \textbf{Описание} \\ \hline
    Код операции & 1 байт & uint8\_t & Уникальный код операции (0x02 (REGISTRATION) — запрос на регистрацию нового пользователя) \\ \hline
    Длина логина & 1 байт & uint8\_t & Количество байт в строке логина \\ \hline
    Длина пароля & 1 байт & uint8\_t & Количество байт в строке пароля \\ \hline
    Логин & L байт (\(L \le 255\)) & byte[L] (UTF-8) & Логин пользователя, длина определяется значением поля "Длина логина" \\ \hline
    Пароль & L байт (\(L \le 255\)) & byte[L] (UTF-8) & Пароль пользователя, длина определяется значением поля "Длина пароля" \\ \hline
    \end{tabular}
    \label{tab:register_request}
\end{table}

\begin{table}[H]
    \centering
    \caption{Атрибутивный состав ответа сервера на регистрацию нового пользователя}
    \renewcommand{\arraystretch}{1.2}
    \begin{tabular}{|p{4cm}|p{2.2cm}|p{3cm}|p{6cm}|}
    \hline
    \textbf{Параметр} & \textbf{Размер} & \textbf{Тип данных} & \textbf{Описание} \\ \hline
    Код операции & 1 байт & uint8\_t & Уникальный код операции (0x02 (REGISTRATION) — запрос на регистрацию нового пользователя) \\ \hline
    Статус выполнения & 1 байт & uint8\_t & 0 — ошибка, 1 — успех \\ \hline
    ID пользователя & 1 байт & uint8\_t & в случае status == 1, передается 0, в остальных случаях нумерация пользователей начинается 0 \\ \hline
    \end{tabular}
    \label{tab:auth_response}
\end{table}



\begin{table}[H]
    \centering
    \caption{Атрибутивный состав запроса на получение списка заметок определенного пользователя}
    \begin{tabular}{|p{4cm}|p{2cm}|p{3cm}|p{6cm}|}
    \hline
    \textbf{Параметр} & \textbf{Размер} & \textbf{Тип данных} & \textbf{Описание} \\ \hline
    Код операции & 1 байт & uint8\_t & Уникальный код операции (0x03 (GET\_NOTES) — запрос на получение списка заметок для определенного пользователя) \\ \hline
    ID пользователя & 1 байт & uint8\_t & Идентификатор пользователя \\ \hline
    \end{tabular}
    \label{tab:note_create_request}
\end{table}


\begin{table}[H]
    \centering
    \caption{Атрибутивный состав ответа сервера на запрос получения списка заметок определенного пользователя}
    \begin{tabular}{|p{4cm}|p{2cm}|p{3cm}|p{6cm}|}
    \hline
    \textbf{Параметр} & \textbf{Размер} & \textbf{Тип данных} & \textbf{Описание} \\ \hline
    Код операции & 1 байт & uint8\_t & Уникальный код операции (0x03 (GET\_NOTES) — запрос на получение списка заметок для определенного пользователя) \\ \hline
    Количество заметок & 1 байт & uint8\_t & Число возвращаемых заметок для определенного пользователя \\ \hline
    ID заметки & 1 байт & uint8\_t & ID очередной возвращаемой заметки \\ \hline
    Версия заметки & 1 байт & uint8\_t & Номер версии очередной заметки \\ \hline
    Флаг указывающий на то совместная заметка или нет & 1 байт & uint8\_t & 0 - заметка личная, 1 - заметка совместная \\ \hline 
    ID владельца заметки & 1 байт & uint8\_t & Идентификатор владельца заметки \\ \hline
    Длина названия заметки & 1 байт & uint8\_t & Длина названия заметки \\ \hline
    Название заметки & L байт (\(L \le 255\)) & byte[L] (UTF-8) & Название заметки, длина определяется значением поля "Длина названия заметки" \\ \hline
    \end{tabular}
    \label{tab:note_create_request}
\end{table}


\begin{table}[H]
    \centering
    \caption{Атрибутивный состав запроса создания новой заметки пользователем}
    \begin{tabular}{|p{4cm}|p{2cm}|p{3cm}|p{6cm}|}
    \hline
    \textbf{Параметр} & \textbf{Размер} & \textbf{Тип данных} & \textbf{Описание} \\ \hline
    Код операции & 1 байт & uint8\_t & Уникальный код операции (0x04 - CREATE\_NOTE) \\ \hline
    ID пользователя & 1 байт & uint8\_t & Идентификатор пользователя \\ \hline
    Размер названия заметки & 1 байт & uint8\_t & Размер названия созданной пользователем заметки \\ \hline
    Название заметки & L байт (\(L \le 255\)) & byte[L] (UTF-8) & Название созданной пользователем заметки, длина определяется значением поля "Размер названия заметки" \\ \hline
    \end{tabular}
    \label{tab:note_create_request}
\end{table}


\begin{table}[H]
    \centering
    \caption{Атрибутивный состав ответа сервера на создание заметки}
    \renewcommand{\arraystretch}{1.2}
    \begin{tabular}{|p{4cm}|p{2cm}|p{3cm}|p{6cm}|}
    \hline
    \textbf{Параметр} & \textbf{Размер} & \textbf{Тип данных} & \textbf{Описание} \\ \hline
    Код операции & 1 байт & uint8\_t & Уникальный код операции (0x04 - CREATE\_NOTE) \\ \hline
    Код результата & 1 байт & uint8\_t & 0 - заметка успешно создана, 1 - заметка с таким именем уже существует для данного пользователя  \\ \hline
    ID созданной заметки & 1 байт & uint8\_t & Идентификатор созданной заметки \\ \hline
    Версия созданной заметки & 1 байт & uint8\_t & Версия созданной заметки \\ \hline
    Длина названия заметки & 1 байт & uint8\_t & Длина названия созданной заметки \\ \hline
    Название заметки & L байт (\(L \le 255\)) & byte[L] (UTF-8) & Название заметки, длина определяется значением поля "Длина названия заметки" \\ \hline
    \end{tabular}
    \label{tab:note_create_response_min}
\end{table}


\begin{table}[H]
    \centering
    \caption{Атрибутивный состав запроса синхронизации заметки}
    \begin{tabular}{|p{4cm}|p{2cm}|p{3cm}|p{6cm}|}
    \hline
    \textbf{Параметр} & \textbf{Размер} & \textbf{Тип данных} & \textbf{Описание} \\ \hline
    Код операции & 1 байт & uint8\_t & Уникальный код операции (0x05 - SYNC) \\ \hline
    ID заметки & 1 байт & uint8\_t & Идентификатор заметки, для которой необходима синхронизация \\ \hline
    ID пользователя & 1 байт & uint8\_t & Идентификатор пользователя, запросившего синхронизацию \\ \hline
    Длина базовой версии текста (до редактирования) & 1 байт & uint8\_t & Длина базовой версии текста (до редактирования) \\ \hline
    Длина локальной версии текста (после редактирования) & 1 байт & uint8\_t & Длина локальной версии текста (после редактирования) \\ \hline
    Текст базовой версии (до редактирования) & L байт (\(L \le 255\)) & byte[L] (UTF-8) & Текст базовой версии (до редактирования), длина определяется значением поля "Длина базовой версии текста (до редактирования) " \\ \hline
    Текст локальной версии (после редактирования) & L байт (\(L \le 255\)) & byte[L] (UTF-8) & Текст локальной версии (после редактирования), длина определяется значением поля "Длина локальной версии текста (после редактирования)" \\ \hline
    \end{tabular}
    \label{tab:note_create_request}
\end{table}


\begin{table}[H]
    \centering
    \caption{Атрибутивный состав ответа сервера на запрос синхронизации заметки}
    \renewcommand{\arraystretch}{1.2}
    \begin{tabular}{|p{4cm}|p{2cm}|p{3cm}|p{6cm}|}
    \hline
    \textbf{Параметр} & \textbf{Размер} & \textbf{Тип данных} & \textbf{Описание} \\ \hline
    Код операции & 1 байт & uint8\_t & Уникальный код операции (0x05 - SYNC) \\ \hline
    Результат синхронизации & 1 байт & uint8\_t & status == 0 - было произведено автоматическое слияние версий, status == 1 - необходимо ручное решение конфликта \\ \hline
    ID заметки & 1 байт & uint8\_t & Идентификатор заметки \\ \hline
    Номер версии & 1 байт & uint8\_t & Если status == 0, то номер версии заметки после автоматического слияния, если status == 1, то номер версии заметки на сервере, \\ \hline
    Длина текста & 1 байт & uint8\_t & Длина текста заметки \\ \hline
    Текст заметки & L байт (\(L \le 255\)) & byte[L] (UTF-8) & Текст заметки, длина определяется значением поля "Длина текста" \\ \hline
    \end{tabular}
    \label{tab:note_create_response_min}
\end{table}


\begin{table}[H]
    \centering
    \caption{Атрибутивный состав запроса открытия заметки}
    \begin{tabular}{|p{4cm}|p{2cm}|p{3cm}|p{6cm}|}
    \hline
    \textbf{Параметр} & \textbf{Размер} & \textbf{Тип данных} & \textbf{Описание} \\ \hline
    Код операции & 1 байт & uint8\_t & Уникальный код операции (0x06 - OPEN\_NOTE) \\ \hline
    ID заметки & 1 байт & uint8\_t & Идентификатор заметки, для которой необходима синхронизация \\ \hline
    ID пользователя & 1 байт & uint8\_t & Идентификатор пользователя, запросившего синхронизацию \\ \hline
    Версия заметки на клиенте & 1 байт & uint8\_t & Версия заметки на сервере \\ \hline
    \end{tabular}
    \label{tab:note_create_request}
\end{table}


\begin{table}[H]
    \centering
    \caption{Атрибутивный состав ответа сервера на запрос открытия заметки}
    \renewcommand{\arraystretch}{1.2}
    \begin{tabular}{|p{4cm}|p{2cm}|p{3cm}|p{6cm}|}
    \hline
    \textbf{Параметр} & \textbf{Размер} & \textbf{Тип данных} & \textbf{Описание} \\ \hline
    Код операции & 1 байт & uint8\_t & Уникальный код операции (0x06 - OPEN\_NOTE) \\ \hline
    Статус & 1 байт & uint8\_t & Результат открытия заметки \\ \hline
    ID заметки & 1 байт & uint8\_t & Идентификатор открытой заметки, в случае если status == 1, то возвращается 0 \\ \hline
    Версия заметки на сервере & 1 байт & uint8\_t & Версия заметки на сервере \\ \hline
    Длина названия заметки & 1 байт & uint8\_t & Длина названия заметки \\ \hline
    Название заметки & L байт (\(L \le 255\)) & byte[L] (UTF-8) & Название заметки, длина определяется значением поля "Длина названия заметки" \\ \hline
    Длина текста & 1 байт & uint8\_t & Длина текста заметки \\ \hline
    Текст заметки & L байт (\(L \le 255\)) & byte[L] (UTF-8) & Текст открытой пользователем заметки, длина определяется значением поля "Длина текста" \\ \hline
    \end{tabular}
    \label{tab:note_create_response_min}
\end{table}



\begin{table}[H]
    \centering
    \caption{Атрибутивный состав запроса обновления текста заметки}
    \begin{tabular}{|p{4cm}|p{2cm}|p{3cm}|p{6cm}|}
    \hline
    \textbf{Параметр} & \textbf{Размер} & \textbf{Тип данных} & \textbf{Описание} \\ \hline
    Код операции & 1 байт & uint8\_t & Уникальный код операции (0x07 - UPDATE\_TEXT) \\ \hline
    ID пользователя & 1 байт & uint8\_t & Индентификатор пользователя запросившего обновление заметки \\ \hline
    ID заметки & 1 байт & uint8\_t & Идентификатор заметки, в которой обновляется текст \\ \hline
    Версия заметки на клиенте & 1 байт & uint8\_t & Версия заметки на клиенте \\ \hline
    Длина обновленного текста заметки & 1 байт & uint8\_t & Обвленный пользователем текст \\ \hline
    \end{tabular}
    \label{tab:note_create_request}
\end{table}

\begin{table}[H]
    \centering
    \caption{Атрибутивный состав ответа сервера на запрос обновления текста заметки}
    \renewcommand{\arraystretch}{1.2}
    \begin{tabular}{|p{4cm}|p{2cm}|p{3cm}|p{6cm}|}
    \hline
    \textbf{Параметр} & \textbf{Размер} & \textbf{Тип данных} & \textbf{Описание} \\ \hline
    Код операции & 1 байт & uint8\_t & Уникальный код операции (0x07 - UPDATE\_TEXT) \\ \hline
    Статус & 1 байт & uint8\_t & Результат открытия заметки \\ \hline
    ID заметки & 1 байт & uint8\_t & Идентификатор обновления заметки, в случае если status == 1, то возвращается 0 \\ \hline
    Версия заметки на сервере & 1 байт & uint8\_t & Версия заметки на сервере после обновления \\ \hline
    \end{tabular}
    \label{tab:note_create_response_min}
\end{table}


\begin{table}[H]
    \centering
    \caption{Атрибутивный состав запроса расшаривания заметки}
    \begin{tabular}{|p{4cm}|p{2cm}|p{3cm}|p{6cm}|}
    \hline
    \textbf{Параметр} & \textbf{Размер} & \textbf{Тип данных} & \textbf{Описание} \\ \hline
    Код операции & 1 байт & uint8\_t & Уникальный код операции (0x08 - SHARE\_NOTE) \\ \hline
    ID пользователя & 1 байт & uint8\_t & Индентификатор пользователя, поделившегося заметкой \\ \hline
    ID заметки & 1 байт & uint8\_t & Идентификатор заметки, которой пользователь решил поделиться \\ \hline
    Версия заметки на клиенте & 1 байт & uint8\_t & Версия заметки на клиенте \\ \hline
    \end{tabular}
    \label{tab:note_create_request}
\end{table}



\begin{table}[H]
    \centering
    \caption{Атрибутивный состав запроса нотификации расшаривания заметки}
    \begin{tabular}{|p{4cm}|p{2cm}|p{3cm}|p{6cm}|}
    \hline
    \textbf{Параметр} & \textbf{Размер} & \textbf{Тип данных} & \textbf{Описание} \\ \hline
    Код операции & 1 байт & uint8\_t & Уникальный код операции (0x09 - SHARE\_NOTE\_NOTIFY) \\ \hline
    ID заметки & 1 байт & uint8\_t & Идентификатор заметки, которой пользователь решил поделиться \\ \hline
    ID владельца заметки & 1 байт & uint8\_t & Индентификатор пользователя, поделившегося заметкой \\ \hline
    Версия заметки на сервере & 1 байт & uint8\_t & Версия заметки на сервере \\ \hline
    Длина названия заметки & 1 байт & uint8\_t & Длина названия заметки \\ \hline
    Название заметки & L байт (\(L \le 255\)) & byte[L] (UTF-8) & Название заметки, длина определяется значением поля "Длина названия заметки" \\ \hline
    \end{tabular}
    \label{tab:note_create_request}
\end{table}





\begin{table}[H]
    \centering
    \caption{Атрибутивный состав запроса нотификации расшаривания заметки}
    \begin{tabular}{|p{4cm}|p{2cm}|p{3cm}|p{6cm}|}
    \hline
    \textbf{Параметр} & \textbf{Размер} & \textbf{Тип данных} & \textbf{Описание} \\ \hline
    Код операции & 1 байт & uint8\_t & Уникальный код операции (0x09 - SHARE\_NOTE\_NOTIFY) \\ \hline
    ID заметки & 1 байт & uint8\_t & Идентификатор заметки, которой пользователь решил поделиться \\ \hline
    ID владельца заметки & 1 байт & uint8\_t & Индентификатор пользователя, поделившегося заметкой \\ \hline
    Версия заметки на сервере & 1 байт & uint8\_t & Версия заметки на сервере \\ \hline
    Длина названия заметки & 1 байт & uint8\_t & Длина названия заметки \\ \hline
    Название заметки & L байт (\(L \le 255\)) & byte[L] (UTF-8) & Название заметки, длина определяется значением поля "Длина названия заметки" \\ \hline
    \end{tabular}
    \label{tab:note_create_request}
\end{table}




\begin{table}[H]
    \centering
    \caption{Атрибутивный состав запроса подтвреждения слияния версий заметки}
    \begin{tabular}{|p{4cm}|p{2cm}|p{3cm}|p{6cm}|}
    \hline
    \textbf{Параметр} & \textbf{Размер} & \textbf{Тип данных} & \textbf{Описание} \\ \hline
    Код операции & 1 байт & uint8\_t & Уникальный код операции (0x10 - APPROVE\_MERGE) \\ \hline
    ID заметки & 1 байт & uint8\_t & Идентификатор заметки, которой пользователь решил поделиться \\ \hline
    ID пользователя & 1 байт & uint8\_t & Индентификатор пользователя, у которого возник конфликт в заметке \\ \hline
    Вариант решения конфликта пользователем & 1 байт & uint8\_t & 0 - ручное решение конфликта, 1 - принятие серверной версии заметки \\ \hline
    Длина текста & 1 байт & uint8\_t & Длина текста заметки, передается только в случае решения конфликта ручным методом \\ \hline
    Текст смерженной версии & L байт (\(L \le 255\)) & byte[L] (UTF-8) & Текст передается только в случае, если пользователь выбрал ручной способ решения конфликта \\ \hline
    \end{tabular}
    \label{tab:note_create_request}
\end{table}



\begin{table}[H]
    \centering
    \caption{Атрибутивный состав запроса подтвреждения слияния владельцем заметки}
    \begin{tabular}{|p{4cm}|p{2cm}|p{3cm}|p{6cm}|}
    \hline
    \textbf{Параметр} & \textbf{Размер} & \textbf{Тип данных} & \textbf{Описание} \\ \hline
    Код операции & 1 байт & uint8\_t & Уникальный код операции (0x11 - OWNER\_APPROVE\_MERGE) \\ \hline
    ID заметки & 1 байт & uint8\_t & Идентификатор заметки, которой пользователь решил поделиться \\ \hline
    ID пользователя & 1 байт & uint8\_t & Индентификатор пользователя, отправившегося запрос на слияние \\ \hline
    Длина текста & 1 байт & uint8\_t & Длина текста заметки после ручного слияния пользователем \\ \hline
    Текст смерженной версии & L байт (\(L \le 255\)) & byte[L] (UTF-8) & Текст смерженной версии заметки, длина определяется значением поля "Длина текста" \\ \hline
    \end{tabular}
    \label{tab:note_create_request}
\end{table}


\begin{table}[H]
    \centering
    \caption{Атрибутивный состав ответа сервера на запрос подтверждения слияния владельцем заметки}
    \renewcommand{\arraystretch}{1.2}
    \begin{tabular}{|p{4cm}|p{2cm}|p{3cm}|p{6cm}|}
    \hline
    \textbf{Параметр} & \textbf{Размер} & \textbf{Тип данных} & \textbf{Описание} \\ \hline
    Код операции & 1 байт & uint8\_t & Уникальный код операции (0x11 - OWNER\_APPROVE\_MERGE) \\ \hline
    Статус & 1 байт & uint8\_t & Результат подтвреждения слияния владельцем пользователя (0 - версия принята, 1 - версия не принята) \\ \hline
    ID пользователя & 1 байт & uint8\_t & Идентификатор пользователя, отправившего версию на согласование \\ \hline
    ID заметки & 1 байт & uint8\_t & Идентификатор заметки \\ \hline
    Версия заметки & 1 байт & uint8\_t & Версия заметки после принятия или отказа владельцем заметки (в случае принятия версия заметки увеличивается, иначе остается той же) \\ \hline
    Длина текста, в версии на слияние & 1 байт & uint8\_t & Длина текста, отправленного на слияние \\ \hline
    Текст заметки, отправленной на согласование & L байт (\(L \le 255\)) & byte[L] (UTF-8) & Текст версии заметки, отправленной на согласование, длина определяется значением поля "Длина текста" \\ \hline
    \end{tabular}
    \label{tab:note_create_response_min}
\end{table}



\begin{table}[H]
    \centering
    \caption{Атрибутивный состав запроса принятия серверной версии}
    \begin{tabular}{|p{4cm}|p{2cm}|p{3cm}|p{6cm}|}
    \hline
    \textbf{Параметр} & \textbf{Размер} & \textbf{Тип данных} & \textbf{Описание} \\ \hline
    Код операции & 1 байт & uint8\_t & Уникальный код операции (0x12 - SERVER\_APPROVE\_MERGE) \\ \hline
    ID заметки & 1 байт & uint8\_t & Идентификатор заметки, которой пользователь решил поделиться \\ \hline
    ID версии & 1 байт & uint8\_t & Идентификатор серверной версии заметки, принятой пользователем \\ \hline
    Длина текста & 1 байт & uint8\_t & Длина текста заметки в результате принятия серверной версии \\ \hline
    Текст смерженной версии & L байт (\(L \le 255\)) & byte[L] (UTF-8) & Текст серверной версии, принятой пользователем, длина определяется значением поля "Длина текста" \\ \hline
    \end{tabular}
    \label{tab:note_create_request}
\end{table}



\begin{table}[H]
    \centering
    \caption{Атрибутивный состав запроса принятой версии в результате слияния}
    \begin{tabular}{|p{4cm}|p{2cm}|p{3cm}|p{6cm}|}
    \hline
    \textbf{Параметр} & \textbf{Размер} & \textbf{Тип данных} & \textbf{Описание} \\ \hline
    Код операции & 1 байт & uint8\_t & Уникальный код операции (0x13 - UPDATE\_TEXT\_MERGED) \\ \hline
    ID заметки & 1 байт & uint8\_t & Идентификатор заметки, которой пользователь решил поделиться \\ \hline
    ID версии & 1 байт & uint8\_t & Идентификатор версии заметки \\ \hline
    Длина текста & 1 байт & uint8\_t & Длина текста заметки после ручного слияния пользователем \\ \hline
    Текст смерженной версии & L байт (\(L \le 255\)) & byte[L] (UTF-8) & Текст смерженной версии заметки, длина определяется значением поля "Длина текста" \\ \hline
    \end{tabular}
    \label{tab:note_create_request}
\end{table}















